\documentclass[
    oneside,                 % stampa solo fronte, evita pagine bianche
    openany,                 % parti/capitoli iniziano sulla prossima pagina
    titlepage,
    11pt,
    a4paper,
    italian,
    DIV=14,                  % specchio di stampa più largo
    headinclude,
    footinclude=true,
    cleardoublepage=empty
]{scrbook}

% --- ClassicThesis + ArsClassica (Pantieri) ---
\usepackage[T1]{fontenc}
\usepackage[utf8]{inputenc}
\usepackage{babel}
\usepackage{lmodern}              % fallback se mathpazo non disponibile per qualcosa
\usepackage{mathpazo}             % Palatino-like per matematica (ClassicThesis default)

\usepackage{amsmath,amssymb,amsthm}
\usepackage{mathtools}
\usepackage{mathrsfs}

\usepackage[
    eulerchapternumbers,
    beramono,
    eulermath,
    pdfspacing,
    parts
]{classicthesis}

\usepackage{arsclassica}          % override Pantieri

% --- nota di colore vinaccia per il numero del capitolo ---
% ArsClassica usa \color{halfgray} per il numero del capitolo;
% ridefinisco halfgray come Maroon (vinaccia ClassicThesis)
\definecolor{halfgray}{rgb}{0.55,0.13,0.16}

% --- aumenta lo spazio dopo il titolo del capitolo per evitare overlap su titoli a 2 righe ---
\usepackage{titlesec}
\titlespacing*{\chapter}{0pt}{2.5\baselineskip}{4\baselineskip}

% --- dimensioni e linguaggio ---
\selectlanguage{italian}

% Aumenta la tolleranza di giustificazione per evitare overfull hbox
\emergencystretch=3em
\tolerance=2000
\hbadness=2000

% --- diagrammi ---
\usepackage{tikz}
\usetikzlibrary{arrows.meta,positioning,decorations.pathmorphing}

% --- grafica e colori (compatibili con ClassicThesis) ---
\usepackage{graphicx}

% --- epigrafi ---
\usepackage{epigraph}
\setlength{\epigraphwidth}{0.55\textwidth}
\renewcommand{\epigraphflush}{flushright}
\renewcommand{\textflush}{flushright}
\renewcommand{\epigraphsize}{\small\itshape}

% --- ambienti teorema-like in stile sobrio ClassicThesis ---
\newtheoremstyle{ctplain}%
  {0.7\baselineskip}%   space above
  {0.5\baselineskip}%   space below
  {\itshape}%           body font
  {}%                   indent
  {\bfseries\scshape}%  head font
  {.}%                  punctuation after head
  {.5em}%               space after head
  {}%                   head spec

\newtheoremstyle{ctdef}%
  {0.7\baselineskip}{0.5\baselineskip}%
  {\normalfont}{}%
  {\bfseries\scshape}{.}%
  {.5em}{}%

\newtheoremstyle{ctrem}%
  {0.7\baselineskip}{0.5\baselineskip}%
  {\normalfont}{}%
  {\itshape}{.}%
  {.5em}{}%

\theoremstyle{ctplain}
\newtheorem{principio}{Principio}[chapter]

\theoremstyle{ctdef}
\newtheorem{defn}[principio]{Definizione}

\theoremstyle{ctrem}
\newtheorem*{oss}{Osservazione}

% --- comandi locali ---
\newcommand{\Lt}{\mathcal{L}_t}
\newcommand{\LP}{\mathcal{L}_P}
\newcommand{\Lo}{\mathcal{L}_o}
\newcommand{\Lr}{\mathcal{L}_r}
\newcommand{\Lf}{\mathcal{L}_f}
\newcommand{\Ld}{\mathcal{L}_d}
\newcommand{\CA}{\mathcal{C}_A}
\newcommand{\CG}{\mathcal{C}_G}
\newcommand{\gA}{g_A}
\newcommand{\gG}{g_G}
\newcommand{\GammaA}{\Gamma_A}
\newcommand{\GammaG}{\Gamma_G}
\newcommand{\Sogno}{\mathfrak{S}}
\newcommand{\selP}{\sigma_P}
\newcommand{\Phii}{\Phi}
\newcommand{\Kappaa}{\mathrm{K}}
\newcommand{\Obs}{\mathrm{Obs}}
\newcommand{\Sym}{\mathrm{Sym}}
\newcommand{\Gs}{G_{\sigma}}
\newcommand{\KSigma}{K(\Sigma^{*})}

% --- hyperref già caricato da classicthesis, configurato in stile sobrio ---
\hypersetup{%
    colorlinks=true,
    linkcolor=Maroon,
    citecolor=Maroon,
    urlcolor=Maroon,
    pdftitle={Bastasse intendersi},
    pdfsubject={note su un'incompatibilità strutturale},
    pdfauthor={}
}

% --- impagina elementi front matter ---
\usepackage{titletoc}

% --- TOC: sentence case ovunque (override classicthesis spacedlowsmallcaps) ---
\renewcommand{\cftpartfont}{\color{CTtitle}}
\renewcommand{\cftpartaftersnum}{}
\renewcommand{\cftpartaftersnumb}{}
\renewcommand{\cftchapfont}{}
\renewcommand{\cftchappresnum}{}
\renewcommand{\cftchapaftersnumb}{}

\begin{document}

\frenchspacing
\raggedbottom
\selectlanguage{italian}
\renewcommand{\partname}{Parte}
\renewcommand{\chaptername}{Capitolo}
\renewcommand{\contentsname}{Indice}
\pagenumbering{roman}
\pagestyle{plain}

% ----------------- frontespizio -----------------
\begin{titlepage}
\begin{center}
\vspace*{2cm}

{\spacedallcaps{note su un'incompatibilità strutturale}}\\[3cm]

{\Huge\bfseries\spacedlowsmallcaps{Bastasse intendersi}}\\[1.5cm]

{\Large\itshape una teoria formale}\\[2cm]

\vfill

{\large\spacedlowsmallcaps{Alvise Vitturi}}\\[0.5cm]

{\itshape Treviso, \the\year}

\end{center}
\end{titlepage}

\clearpage

% ----------------- nota al lettore -----------------
\chapter*{Nota al lettore}
\addcontentsline{toc}{chapter}{Nota al lettore}

Quattro cose, prima di cominciare.

\paragraph{Origine.} Il testo nasce da una relazione vera. È scritto vent'anni dopo, quando il caso è chiuso da entrambi i lati e non si può più rivisitare. Ha una funzione precisa per chi lo scrive: dare struttura a un'osservazione fatta allora, verificata adesso.

\paragraph{Assunzioni.} Sono dichiarate ed esplicite nella Parte~II, in linea di principio falsificabili. Il documento distingue ovunque tra \emph{tesi strutturale generale} --- che vale per la classe --- e \emph{applicazione al caso} --- che è un'interpretazione, e potrebbe essere sbagliata. Il caso non viene anonimizzato. L'anonimato simulerebbe un'oggettività che non c'è: chi scrive è dentro al caso, non sopra di esso.

\paragraph{Forma.} Dove la prosa basta, c'è prosa. Dove la prosa rende sfuocato qualcosa che chiede precisione, entra il formalismo, dichiarato come strumento e non come dimostrazione. I principi numerati sono \emph{enunciati assiologici} --- cose che si assumono per vedere cosa segue --- presentati nella forma del teorema perché quella forma costringe a essere precisi su cosa si sta dicendo, non perché qualcosa sia stato dimostrato.

\paragraph{Architettura.} Una persona è una sezione di un fascio di funtori sopra una categoria di contesti. Le leggi che la attraversano --- $\Lt$, $\LP$, e altre --- sono quei funtori. Non si derivano da una macchina sottostante: affiorano dalla materia che la Parte~I attraversa, e la Parte~II ne dà la forma precisa nello schema rottura-di-simmetria $+$ Noether. Da qui un asse algebrico --- la coordinata di rango e torsione del codominio della carica --- su cui le leggi si dispongono come punti di uno spettro. Tutto il resto del documento --- tempo, metacognizione, errore di tipo, posizionalità linguistica, incompatibilità --- è geometria di questo spettro.

Lo scopo non è convincere nessuno. Non è accusare nessuno. Non è prevedere niente. È fissare. Quello che fissa, resta. Quello che non fissa, è fuori.

\bigskip

\noindent\textbf{Impianto filosofico.} Sei ancoraggi orientano la scrittura. \emph{Il tipo come costruzione} (Martin-L\"of): un giudizio non formabile in una legge non è negato, è non costruibile in quel contesto. \emph{Il tempo come emergenza} (Rovelli): non c'è meta-tempo neutrale, ogni tempo è prodotto dalle interazioni che lo generano. \emph{L'agente come effetto}: non c'è un soggetto a cui le leggi si applicano, c'è una famiglia di leggi da cui il soggetto si lascia leggere. \emph{La necessità senza colpa} (Weil): l'ananke attraversa, e l'attenzione che vede è attenzione che decostruisce l'io che impone. \emph{La forma dell'invarianza} (Landau): la legge e la sua simmetria si nominano insieme; la carica conservata è il contenuto della legge, la simmetria è la forma del suo essere invariante. \emph{La marea, non il martello} (Grothendieck): il formalismo non rompe la materia, la circonda finché la sua forma si lascia formulare. \emph{L'oggetto è il pattern dei suoi sondaggi} (Yoneda): la legge non è entità in sé che il formalismo descrive --- è il pattern delle risposte ai sondaggi attraverso cui si lascia leggere. Le leggi affiorano dalla materia; il formalismo --- monoidi, fasci, cariche conservate --- nomina con precisione la forma che hanno, non la macchina da cui si deriverebbero.

\bigskip

\noindent\textbf{Nota storica.} Il modello è stato riformulato più volte. Le stesure precedenti descrivevano l'incompatibilità prima in termini di due persone con leggi opposte, poi come scelta tra arricchimenti monoidali e quoziente. La stesura attuale completa l'inversione finale: le leggi non sono di persone, sono funtori sopra una categoria di contesti, e una persona è una sezione del fascio che esse formano. L'incompatibilità tra $\Lt$ e $\LP$ non è scelta editoriale di metterle in categorie diverse: è massima distanza spettrale sull'asse algebrico del rango. Tutto il resto del documento --- tempi, metacognizione, errore di tipo, posizionalità --- discende da questa unica coordinata.

\clearpage

% ----------------- indice -----------------
{\setlength{\parskip}{0pt}\tableofcontents}

\clearpage
\pagenumbering{arabic}
\pagestyle{scrheadings}

% ----------------- PARTE I -----------------
\part{Il caso}

\chapter{Prima volta, due anni prima}

{\spacedlowsmallcaps{Avevo quattordici anni.}} Ci eravamo messi insieme --- lei a Vicenza, io a Mestre. L'unico punto di aggancio praticabile era Asiago, e dipendeva da un amico. Ci siamo visti poco.

Avevo notato qualcosa di strano, ma in modo epidermico, e non ci diedi peso. In auto, dietro, con la madre --- Donatella, la madre di Tommaso --- lei era troppo ferma, rigida. Io guardavo fuori dal finestrino. Frecciate fuori contesto, pensate per ferire senza che ci fosse motivo. L'episodio della sorella: la baciai a stampo il giorno dopo che ci eravamo messi insieme; uscita la sorella mi dice \emph{perché mi hai baciato}; le risposi che era una carezza, che non volevo metterla in imbarazzo davanti alla sorella, nessun problema. Andai a Vicenza; dopo una giornata passata insieme a Tommaso, le diedi un bacio e lo trovai difficoltoso. Mi dissi: vabbè, sarà stata arrabbiata.

Per messaggio, dopo, ero distante. Avevo quattordici anni, non ero perso, e la cosa non funzionava --- distante è quello che ero davvero. Un giorno mi arrivò un sms: qualcosa nella forma \emph{ti lascio ma non vorrei, mi hai obbligato} --- non ho il verbatim, ho la firma. Non era diretto. Era una formula con dentro un'attribuzione --- come se la decisione fosse sua ma la causa mia, qualcosa che io le avevo fatto fare. \emph{Mi hai obbligato a cosa?} Non c'era un cosa. Risposi va bene. Il motivo dell'ok fu semplice: non ci aveva messo pathos nessuno dei due, nemmeno lei. Accettai che non avesse funzionato. Per due anni, nient'altro.

\chapter{Il magone alla porta}

Due anni dopo, qui comincia la storia che il documento descrive.

Lei mi mancava --- non sapevo cosa fosse. Le scrissi che sarei andato in montagna. Per messaggio misi il tema sul tavolo: noi. Volevo capire se potevamo riprovarci. Non era ambiguo. Era esplicito, dichiarato. Una sequenza di messaggi in cui io stavo dicendo, in chiaro, quello che volevo --- non sottinteso, non implicito.

Lei rispondeva, ma non rispondeva al tema. Io le scrivevo del noi. Lei mi rispondeva sul colore degli occhi. \emph{Di che colore ho gli occhi?} Verdi, risposi. \emph{Sei stato fortunato.} Non risposi niente. Pensai: fortunato? Se uno se lo ricorda, non è fortuna. Ma soprattutto pensai: gli occhi --- adesso? Stiamo parlando di noi, e tu mi chiedi gli occhi.

Quello che lei non scrisse mai, in nessuno di quegli scambi, fu una qualunque menzione di un ragazzo. Mai un \emph{guarda, prima di proseguire devo dirti una cosa}. Mai un \emph{senti, c'è un contesto che non sai}. Il tema era sul tavolo, esplicitato da me. La risposta non c'era. C'era altro --- domande sul colore degli occhi, parole come \emph{fortunato}. L'informazione che avrebbe chiuso la conversazione, o l'avrebbe spostata su un altro piano, non veniva data.

Arrivai ad Asiago con delle aspettative --- aspettative coerenti con quello che avevo messo per scritto e con il fatto che la conversazione era continuata. Quei condomini orizzontali di montagna dove c'è un giardino comune e le porte al primo piano si affacciano su un marciapiedino, così che esci da una porta ed entri nell'altra. Donatella mi chiese di portare di là qualcosa. Le suonai.

Quando aprii la porta, mi venne un magone. Era bella, per me, oltre ogni aspettativa. La guardai negli occhi, abbassai lo sguardo, lo rialzai: volevo farle passare che la cosa, per me, era stata una sorpresa. Che qualcosa c'era ancora.

Lei lo vide. Era impossibile non vederlo. E nemmeno in quel momento, con il magone passato tra noi sulla soglia, mi disse che c'era un ragazzo.

\begin{oss}[Sulla natura del magone]
È un punto che vale la pena nominare per quello che è, perché altrimenti la prosa lo fa scivolare via come un dato emotivo passeggero. Quello che mi prese sulla soglia non era infatuazione, non era capriccio, non era un'idea che mi ero fatto io e che avrei potuto disfare. Era il \emph{trasalimento} che la tradizione stilnovista aveva già nominato con precisione sette secoli prima, in Cavalcanti: \emph{chi è questa che vèn, ch'ogn'om la mira, / che fa tremar di chiaritate l'âre, / e mena seco Amor, sì che parlare / null'omo pote, ma ciascun sospira?} L'apparire dell'altra che fa tremare l'aria, che impone silenzio, e che porta Amore con sé non come oggetto di scelta ma come dato che attraversa chi guarda. Il movimento dello sguardo --- giù, su --- non era strategia comunicativa: era esattamente l'incapacità di parlare di cui Cavalcanti scrive, sostituita dal piccolo segno che restava praticabile.

Nominarlo così non è ornamento letterario. È il modo più preciso per dire che quel momento appartiene alla stessa famiglia strutturale dell'ananke che la Parte~III nominerà: non lo si decide, lo si riceve. La fenomenologia del trasalimento --- esattamente l'``apertura senza corazza'' che la Parte~III chiamerà Patroclo --- è la postura sotto cui $A$ è arrivato alla soglia, e sotto cui sarebbe tornato al jazz bar. Quella postura è già il dato; tutto il resto del documento articola cosa è successo dopo che quella postura ha incontrato $\LP$.
\end{oss}

\chapter[Il jazz bar]{Il jazz bar: la disequazione e la crepa}

Dopo gli SMS espliciti, dopo il magone passato sulla soglia con lei che lo aveva visto, lei propose una cosa: andiamo in un bar nuovo, suonano jazz. La proposta arrivò tramite Tommaso ma era sua. Una sera in più, in un contesto diverso da casa, con musica e penombra. Costruita conoscendo il quadro: quello che le avevo scritto, quello che le avevo passato sulla porta.

Andammo.

Nel locale, l'atmosfera, lei. Mi dissi: ma dio benedetto, devo proprio farmi problemi. È una bella ragazza, sono cotto, lo ammetto, vado. Mi alzai e dissi una cosa semplice. \emph{Greta, so che due anni fa le cose non sono andate. Non ti chiedo di fidarti di me. Ti dico che io vorrei costruire con te quella fiducia. Ci voglio provare veramente.}

Lei partì con un pianto oltre ogni misura.

Mi avvicinai per capire se potevo aiutarla. In parte ero anche felice --- era evidente che non le ero indifferente. Mi inginocchiai, perché lei era seduta sulla poltrona; non come un cavaliere, per portare il viso alla sua altezza. Vidi gli occhi sbarrati. \emph{Posso andare in bagno?} Certo. Stette in bagno quindici minuti.

Uscì, e non mi cagò. Algida. Andò da Tommaso: \emph{andiamo}. Lì, dentro, mi sono detto: ma che cazzo succede.

Tornammo verso casa. Lei parlava di lingue con Tommaso --- inglese, cosa le piaceva. Lui era imbarazzato. Lei era ilare. Cinque minuti prima il pianto, gli occhi sbarrati, il bagno per quindici minuti, l'uscita algida; adesso ilare, leggera, in conversazione. Io dietro, con le mani in tasca, guardavo le stelle. So distinguere i pianeti. Affioravano Saturno e Giove tra i pini. E mi chiedevo: ma che cavolo. Non ha senso. Questa distanza --- doveva rispondermi anche con un no, e invece niente.

Arrivammo a casa, le due porte adiacenti. Le dissi: \emph{ok Greta, immagino già la risposta, ma è giusto che tu me la dia. Possiamo riprovarci?} Mi guardò. \emph{Non posso. Vorrei, ma ho il ragazzo.}

Lì, dentro, si è aperto qualcosa che prima non c'era. Non rabbia soltanto --- una specie di chiarezza fredda. \emph{Adesso} me lo dici. Non negli SMS, non sulla porta, non prima di andare al bar. Adesso.

Ti chiederò un bacio, pensai, perché voglio capire quel pianto cosa cavolo fosse. \emph{Ok Greta. Ti chiedo solo un bacio per chiudere la storia che abbiamo chiuso solo con un sms.} Sembrò pensarci. Si avvicinò e mi diede un succhiotto sul collo. Un bacio sexy --- decidi tu come chiamarlo.

Pensai: ti ho chiesto un bacio tenero. Cosa cavolo c'entra. E pensai anche: cappero, non pensavo fosse così brava. Molto più brava di me su questo. La distanza tra quello che avevo chiesto e quello che avevo ricevuto era esatta come una riga.

\begin{oss}[La disequazione]
La sera dopo arrivò F, il suo ragazzo. Grigliata, noiosissima. Loro due erano due sconosciuti. Zero baci. Zero abbracci. Zero di nulla. Noia totale.

E in quella noia ho capito una cosa che non avrei mai immaginato di dover capire, perché sembra falsa: F era preferito a me. Non ``preferito moralmente'' o ``scelto perché più serio'' o qualunque cosa avrei potuto raccontarmi. Preferito \emph{come posizione}. F era il ragazzo, io ero qualcun altro. La scena della porta, il pianto, il succhiotto, l'invito --- niente di tutto questo spostava F dalla sua posizione. Era ferma nonostante tutto.

In una formula:
\[
(A,G) \not< (F,G) \quad \text{e però} \quad (F,G) = \text{due morti che camminano.}
\]
E qui c'era qualcosa che non quadrava con nessuna spiegazione che mi fosse familiare. La parola ``preferire'' stava facendo un lavoro diverso da quello che io le chiedevo di fare. Non sapevo ancora quale.
\end{oss}

\begin{oss}[Il timing]
L'informazione sul ragazzo esisteva da prima degli SMS. Da prima del bar. Da prima del pianto. Non è che è mancato il momento per dirla. Sono stati saltati, in ordine: gli SMS in cui avevo messo il tema; la porta in cui lei aveva visto il magone; la sera al bar che lei stessa aveva proposto; il momento in cui mi alzavo e parlavo. È uscita solo dopo --- dopo il pianto, dopo il bagno, dopo l'ilarità, sulla soglia.

Il timing non è inerzia. Non è timidezza, non è ``non sapere come dirlo''. È funzionale: l'informazione esce nel punto in cui può essere usata --- come chiusura di una sera che era stata costruita dopo la mia dichiarazione esplicita. Detta prima, avrebbe chiuso gli SMS o spostato la sera su un altro piano. Detta dopo, lascia che la sera abbia avuto luogo --- pianto compreso, bacio compreso --- e poi sigilla.

La disequazione mi diceva che lei usava una metrica che io non capivo. Il timing mi diceva che usava un calendario che io non capivo. La firma era la stessa: ordinava le persone e distribuiva le informazioni per qualcos'altro --- non per quello che la persona davanti era, non per quando l'informazione sarebbe servita alla persona davanti.

Trovare il nome di quel qualcos'altro era tutto il problema.
\end{oss}

\bigskip

\noindent\textbf{Sull'ordine genetico.} Questa osservazione, allora, era il dato. La sua spiegazione strutturale è in Parte~II e Parte~III, costruita molti anni dopo. Ci sono voluti vent'anni per nominare con precisione cosa rendeva non assurda quell'osservazione semplice.

\chapter[Il giardino]{Il giardino, l'ammorbidimento, la madre}

Il giorno dopo, di mattina, Greta arrivò. C'eravamo io, Tommaso, Donatella. \emph{Vieni in centro, anche Tommaso.} Risposi: \emph{no, scusa dopo ieri sera. Stante le cose, non ho voglia. Se hai cambiato idea, torna quando ti sei chiarita.}

Il pomeriggio, in giardino, mi raggiunse. Mi abbracciò. Mi raccontò di sé --- di quanto le ero mancato in quei due anni, di amiche, di trama. Trama abbondante, e --- ne ero sicuro --- inventata. Era appoggiata, niente di più. Sua madre, dal balcone, la richiamò con tono minaccioso. Mi chiesi: per quale motivo una madre dovrebbe intervenire per qualcosa che non sta accadendo? Eravamo in giardino, in pieno giorno, con un abbraccio tra due ragazzi di sedici anni. Sapessi che immoralità.

Avvicinamento dopo chiusura, senza che la chiusura sia stata revocata. La madre interrompendo, prima che l'avvicinamento producesse qualcosa che andasse onorato. Una mossa, e la sua chiusura preventiva. Tornai a casa con la sensazione che avevo già iniziato ad avere al bar: la scena gira, gli argomenti cambiano.

\chapter{Tommaso al parco: la minaccia}

Per un po' la cosa continuò così. Lei mi cercava, io ero gentile. Può anche essere che qualche volta le abbia scritto io --- ma dopo misi in chiaro che non ero lì per amicizia.

Poco dopo, scalation. Tommaso mi raggiunse al parco, a Mestre. \emph{Greta dice che F vi ha scoperti, e se non la smetti viene a picchiarti.}

Risposi: primo, scoperti di cosa. Il tema è di una banalità semplicissima: io mi sono solo rifiutato di ritrarre il mio sentimento. Se lei gioca sull'ambiguità, ne ho le palle piene anch'io. Senti: io non le scrivo più. Diglielo. Se F vuole parlare con me, non ho l'auto, ma può venire quando vuole.

Poi non ci sentimmo più.

\begin{oss}[La minaccia come copione]
La minaccia non passa direttamente da F. Passa attraverso lei e Tommaso. F è l'\emph{argomento} del messaggio, non il mittente. Il copione assegna a F il ruolo di minaccia per la durata del messaggio --- perché in quella configurazione serve quel ruolo --- senza che F lo abbia richiesto e probabilmente senza che ne sappia niente. È lo stesso meccanismo per cui prima a F era stato assegnato il ruolo di ragazzo, durante la sera al bar, mentre la sera prima l'ordine era diverso. I ruoli sono operatori del copione, e gli argomenti sono intercambiabili.
\end{oss}

\chapter{Vicenza, sei mesi dopo: l'ariete}

Sei mesi dopo, Nicolò (N) mi disse che andavamo a Vicenza. Greta ci aveva invitati per un aperitivo.

Lì venne fuori che adesso aveva T. Lo descriveva lei, e mentre lo descriveva io stavo capendo qualcosa che al jazz bar, sei mesi prima, avevo solo intuito. Lei chiamava T \emph{il quadrupede} --- non con quella parola, ma con la sostanza. Uno con cui aveva avuto una crisi e che, per recuperarla, le aveva portato dei fiori. Era questo che mi raccontava: il quadrupede che porta i fiori.

E mentre lei raccontava T, ho fatto la somma. Sei mesi prima F era preferito a me. Adesso T era preferito a me. La struttura non era cambiata. L'argomento sì.

Il pomeriggio fu un suo monologo. T che ha fatto, T che ha detto, T che dopo la crisi era venuto con i fiori. Io ascoltavo. A un certo punto, dentro la stessa stanza in cui aveva appena finito di descrivermi T, disse: \emph{ho casa libera}.

Lo dici a me. Adesso, dopo questo pomeriggio.

Avevo già fatto tutti i check logici. Il dubbio non è mai esistito. Mi sono detto: uso un bell'ariete. E dicendolo, avrei detto l'unica cosa vera che si poteva dire da seduto in quella stanza, dopo quel monologo: l'unico desiderabile qui è quello che non è in classifica.

\emph{Casa libera? Uhm, ok. Bene, mi sembra un'ottima occasione per fare sesso. Cosa dici?}

Mi guardò, sorrise, rispose di no.

Mi voltai verso Nicolò. \emph{Andiamo a prendere il treno.} Io e lui davanti, lei dietro. Lei disse: \emph{e comunque è anche tardi.} Mi fermai, mi voltai. \emph{Per me nessun problema, i biglietti devo ancora prenderli.} Lei restò in silenzio. Ripresi a camminare verso la stazione.

Sapevo di aver buttato una bomba nucleare per uscire. Però ero stufo.

\begin{oss}[Funzione dell'ariete]
\emph{Casa libera}, in quel pomeriggio, non era proposta. Era la stessa struttura del succhiotto al bar: alza la temperatura, registra la reazione, non fissa conseguenze. L'unica risposta che lascia il copione attivo è una che si presta a essere rimaneggiata --- imbarazzo, esitazione, richiesta di chiarimento. Prendere alla lettera --- dire, piano e secco, \emph{``bene, occasione per fare sesso, cosa dici?''} --- non si presta a essere rimaneggiato. E infatti la risposta è no, sorridendo, e poi \emph{è tardi}: un tentativo di non far chiudere la scena.

C'è però una seconda funzione, e per me contava di più. Perché arrivarci, va escluso uno scivolamento facile: la spiegazione comoda sotto cui questa storia si liquiderebbe in un paragrafo, ovvero \emph{il cazzone duro vince e il sentimentale perde}. Quella spiegazione qui non regge, e non per ragioni mie --- per ragioni sue. Nel modo in cui lei stessa raccontava F e T, né l'uno né l'altro era distaccato o duro. T è quello che dopo la crisi viene con i fiori. F è quello con cui non c'è neanche un abbraccio per assenza di postura, non per durezza. Erano ragazzi cedevoli, dipendenti, esecutivi --- anche per come li descriveva lei. Non vinceva chi era duro. Vinceva qualcos'altro, su un asse a cui né la mia esposizione né l'ipotesi di durezza simulata avrebbero cambiato la posizione.

Allora l'ariete non era solo uscita dall'ambiguità. Era uscita dall'intero sistema di assi che mi era familiare. Non potevo dichiararmi di nuovo --- l'avevo già fatto al bar. Non potevo simulare durezza --- la durezza nemmeno era quello che vinceva. Ogni mossa che restava dentro un asse a me leggibile era una mossa già sconfitta. L'unica mossa sensata era una che uscisse dagli assi. Pronunciandola, ero già fuori. Prima ancora che lei rispondesse no.
\end{oss}

\chapter{Canottaggio}

Tra Vicenza e il compleanno c'è un episodio che vale la pena raccontare per quello che mostra, non per quello che pesa. Quel periodo lo trascorrevo facendo canottaggio. N veniva spesso. Un giorno, in doccia, mi disse \emph{ehi, guarda cosa mi ha mandato Greta}. Tirò fuori il telefono. C'era una sua foto, nuda.

Restai zitto. Quello che vedevo, in ordine non logico ma simultaneo, era questo: lei stava tradendo T \emph{anche} con N, e probabilmente non solo; lei lo aveva dichiarato esplicitamente con una foto, non con un'allusione; N era uno a cui mostrare quella foto in doccia a un altro --- a me --- sembrava una cosa normale da fare. Le tre osservazioni stavano insieme. Non potevo separarle.

\begin{oss}[Cosa Canottaggio mette in chiaro]
\emph{La configurazione è esposta.} La foto a $N$ non era dichiarazione cifrata. Era la configurazione del pool, esposta come $\kappa(P)$: $T$ era in posizione bassa, $N$ in posizione alta, l'esecuzione era in corso. Chiunque vedesse la foto vedeva anche la configurazione. Non c'era niente da decifrare.

\emph{Il pianificatore non c'è.} Un pianificatore strategico --- la figura che $\Lt$ proietta dentro $\LP$ --- non manda quella foto a $N$. Sa che $N$ è uno che la girerà, sa che $N$ è amico mio, sa che quella foto può arrivare al compleanno preparata. Quel pianificatore tiene fuori $N$ dalla cosa, oppure manda una foto meno netta, oppure dispone le pedine perché la configurazione non viaggi. Niente di tutto questo accade, perché il copione non computa effetti collaterali a valle del proprio passo corrente. Il passo corrente era ``segnale alto-temperatura verso $N$''. La conseguenza --- foto in doccia a $A$ --- è fuori dalla legge.

\emph{La scusa ``ho il ragazzo'' era morfismo locale.} Quella frase, detta nei giorni del jazz bar e di Vicenza per chiudere un mio passaggio, era stata letta da me come dichiarazione persistente. La foto a $N$ la contraddice senza rimedio. Non è bugia, e non era cortesia: era esecuzione di un morfismo richiesto in quella configurazione, che non si impegna oltre il proprio passo. Sotto $\Lt$ è doppiezza. Sotto $\LP$ è regolarità.

\emph{Il giardino del giorno dopo.} Letto da $\Lt$, il mattino dopo il jazz bar è incomprensibile: dopo una sera con quel peso, ``vieni in centro, anche Tommaso'' non avrebbe forma né di riavvicinamento né di distanza. Era come se la sera prima non fosse accaduta nel suo tempo proprio. Era così, infatti: la sera prima era esecuzione chiusa di un copione, la mattina era nuova esecuzione su una configurazione nuova. La foto a $N$ è la stessa firma a distanza di mesi.

\emph{N non era $\Lt$.} N stava con Greta esattamente nel registro che il suo copione produceva: episodi, scambi, niente da custodire. Mostrare la foto in doccia a un amico non era violazione di un'intimità --- non c'era un'intimità persistente tra lui e lei da custodire, dal suo punto di vista. Era nello stesso tempo locale, e nello stesso tempo locale quel gesto era congruente. La struttura ha trovato in $N$ un argomento senza attrito --- non perché $N$ fosse stupido o cattivo, ma perché abitava il registro temporale giusto per quel ruolo.

\emph{Nemmeno strategia di ingelosimento.} La lettura più tenera che $\Lt$ può ancora provare a fare è: lei la mandava a $N$ \emph{per me}, sapendo che $N$ me l'avrebbe mostrata, calcolando l'effetto su di me. È una lettura che non regge nemmeno dentro $\Lt$. Una strategia di ingelosimento richiede un canale di trasmissione del segnale: io e lei non ci sentivamo, lei non mi cercava, io neppure. Richiede controllo sulla consegna: la foto poteva non arrivarmi mai. Richiede verifica dell'effetto: lei non aveva modo di sapere se l'avevo ricevuta, né se avevo reagito. Senza canale, senza controllo, senza verifica, la ``strategia di ingelosimento'' è un'ipotesi che cade da sola. La foto è arrivata in doccia a me per pura propagazione del pool --- $N$ era amico mio, $N$ ha mostrato. Non per design. Supporre che fosse per me, in assenza di tutto quello che renderebbe la supposizione plausibile, è egocentrismo. È il falso testimone (Parte~II) al suo stadio più sottile: non solo $\Lt$ inventa un pianificatore dove non c'è, ma gli attribuisce \emph{me} come destinatario delle sue mosse. Il pianificatore-per-me è proiezione al quadrato. Il fatto strutturale è più piano: il copione gira sui propri argomenti correnti, e gli effetti collaterali si propagano nel pool senza essere stati previsti.
\end{oss}

A Canottaggio, in doccia, davanti al telefono di N, sapevo già tutto quello che il compleanno avrebbe esposto. La configurazione era pubblica da settimane. Quando settimane dopo, al compleanno, tradì $T$ con $N$ ``in modo lampante, per dimostrazione'', non era novità --- era esecuzione coerente di una configurazione che avevo già letto. La mia non-reazione al compleanno non era forza d'animo. Era assenza di informazione nuova.

\chapter{Il compleanno: il tradimento}

Andai un'altra volta in montagna. Poi mi invitò al suo compleanno.

A quel compleanno tradì T con N --- facendomelo capire in modo lampante. Non per inavvertenza. Per dimostrazione. Perché io vedessi che era in grado di farlo, perché io vedessi che N era stato scelto per quella funzione, perché io vedessi.

Non ho reagito. Mi sono mangiato la pizza, ho festeggiato il compleanno, e poi sono tornato da Sara. Mi ero già trovato una ragazza da un pezzo, ma lei continuava a non integrare la cosa.

\begin{oss}[Il compleanno come configurazione esposta, non come messaggio]
Sotto $\LP$, il compleanno con T e N non è scena costruita per A. $\LP$ non costruisce scene per destinatari --- non ha la grammatica per farlo. È una configurazione del pool: T accanto ($\mathrm{autonomy} \approx 0$), N in alto in classifica per accumulazione di $\Sogno$ e decadimento di $\mathrm{threat}$ (supply disponibile), la festa come campo. La configurazione si è prodotta.

Letta da A in $\CA$, la stessa configurazione produce un'altra cosa. La metrica di $\Lt$ non vede la configurazione come configurazione: la vede come scena. \emph{Guarda di cosa sono capace, guarda chi mi sta accanto adesso, guarda cosa avresti potuto avere se fossi stato il tipo giusto}. La lettura è interna ad A, non è proprietà del fatto. Le tre frasi sono quello che la geometria di $\CA$ produce quando il pool di lei si configura così; non sono quello che lei stava dicendo, perché lei, sotto $\LP$, non stava dicendo. La configurazione, su un osservatore in $\CA$, si traduce nelle frasi che A si ritrova in mente. Nessuna delle due letture è falsa; sono due letture diverse dello stesso fatto da due leggi diverse.

$\LP$, da parte sua, espone. Non per comunicare --- $\LP$ non vuole nulla --- ma perché $\LP$ non ha altro modo di esistere che mostrare la propria configurazione sulla base. Le manca il canale, le resta la presenza. Il compleanno con T e N è la presenza di $\LP$ esposta al pool che la circonda, A incluso. Lo stato di lei era quello, e quello stato si vedeva.

La configurazione non era nuova. Canottaggio l'aveva esposta settimane prima --- la foto a N in doccia era già la stessa configurazione. La non-reazione di A al compleanno --- pizza, festeggiamenti, ritorno da Sara --- non è equanimità conquistata. È che la configurazione era già stata letta. Una manifestazione, per essere registrata come notizia in $\CA$, ha bisogno di non essere già nota; arrivava su un osservatore che aveva già visto.

La storia finisce qui non perché A abbia trovato la risposta giusta, ma perché la configurazione che il compleanno mostrava era stata letta abbastanza volte da non costituire più atto in $\CA$. $\LP$ non aveva fatto altro che continuare a essere $\LP$; A non aveva fatto altro che finire di leggere.
\end{oss}

\chapter{Cosa chiede una struttura}

Questi episodi --- la disequazione del bar, l'ammorbidimento del giardino, la minaccia per interposta persona, l'ariete a Vicenza, la dimostrazione del compleanno --- chiedono una struttura per essere capiti. Le narrazioni disponibili non bastano.

Non era cinismo: il pianto al bar era reale, il magone alla porta era reale. Non era malvagità, non come dato positivo: lei non si sentiva malvagia mentre lo faceva, e il giudizio morale ordinario qui descrive male. Non era immaturità: la sequenza ha durato anni, con coerenza interna, in più contesti, su più destinatari (F, T, e me). Non era errore: gli errori si correggono dopo la prima evidenza, e qui la correzione non c'era. Non era confusione: l'informazione decisiva --- l'esistenza del ragazzo --- è stata trattenuta esattamente fino al punto in cui la sera era stata già spesa, e il timing di una mossa così precisa non si produce per disorientamento. E non era nemmeno la spiegazione classica --- non era ``preferire il duro al sentimentale''. F e T, anche per come li descriveva lei, non erano duri. Non c'era un asse di mascolinità sotto cui lei vinceva e io perdevo.

Quello che spiega la sequenza è qualcos'altro --- più strano e più strutturale. Non un difetto da rimuovere. Una geometria interna in cui quegli atti seguivano la propria coerenza, mentre nella mia geometria ne seguivano un'altra. Le due geometrie non si parlavano. Ogni passaggio dell'una verso l'altra costava qualcosa, e ad alcune distanze il costo diventava la cosa.

Una nota da fissare prima della Parte~II, per non lasciare un buco implicito. Ciò che rendeva $F$, $T$, $N$ argomenti del copione --- ciò che li portava in cima al pool nei momenti in cui ci salivano --- non era una proprietà che la legge poteva calcolare. Era un orientamento pre-riflessivo della persona che sognava: tropismi, attrazioni, gradi di intensità con cui quegli argomenti, per lei, erano vivi. La legge $\LP$ filtrava quell'orientamento --- escludeva chi portava troppa autonomia, premiava chi non disturbava la mappa --- ma non lo produceva. Il sogno era suo, e restava suo. Il documento, sulla forma del filtro, può dire; sul contenuto dell'orientamento, no --- e dichiarerà questa opacità come parte del bordo del logos.

La Parte~II prova a dare un nome a quella geometria. Senza pretesa di dimostrazione --- come strumento per nominare con precisione qualcosa che la prosa, da sola, rende sfuocato.

% ----------------- PARTE II -----------------
\part{La forma}

\noindent\itshape Una sola idea, percorsa in cinque capitoli. Una persona è una sezione di un fascio di funtori sopra una categoria di contesti; le leggi che la attraversano sono quei funtori; tutto il resto --- tempo, metacognizione, errore di tipo, posizionalità, incompatibilità --- è geometria di questa sezione. Il formalismo è strumento, non descrizione ontologica: nomina la forma del pattern, non cosa una persona è. I principi numerati sono enunciati assiologici, non risultati derivati.\upshape

\bigskip

\chapter{La materia e i suoi sondaggi}

Un contesto non è il mondo. È un posto strutturale in una persona --- \emph{Greta-prima-del-jazz-bar}, \emph{il padre}, \emph{Vicenza-quel-pomeriggio} --- nominato dalla persona stessa, mai dall'esterno. Una persona contiene molti contesti, e i contesti si parlano: uno si trasforma in un altro (\emph{Greta} $\to$ \emph{ex-Greta}); uno deriva da un altro (\emph{amici-università} $\to$ \emph{Caio}); uno riconosce in un altro la stessa forma (\emph{Greta} $\to$ \emph{Veronica}, come ripetizione strutturale).

\begin{defn}[Categoria dei contesti]
$\mathbf{C}$ è la categoria che ha per oggetti i contesti relazionali attivi di una persona, e per morfismi tre famiglie: \emph{evoluzione} ($c \to c'$, $c$ è diventato $c'$, irreversibile), \emph{derivazione} ($c \to c'$, $c$ è genitore di $c'$), \emph{affinità} ($c \to c'$, $c$ e $c'$ sono riconosciuti dalla persona come la stessa forma).
\end{defn}

$\mathbf{C}$ non è canonica. È la $\mathbf{C}$ che la persona costruisce di sé. Il documento che segue è la $\mathbf{C}$ dell'autore vent'anni dopo --- con tutta la riscrittura retrospettiva che il tempo permette. La soggettività della costruzione va dichiarata: nessuna $\mathbf{C}$ neutra esiste.

\section*{La legge come funtore}

Una legge $L$ è uno schema $(S_L, \Gamma_L, \rho_L)$ secondo Noether: $S_L$ è una struttura algebrica, $\Gamma_L$ il suo gruppo di simmetria pieno, $\rho_L$ la rottura $H_L \subset \Gamma_L$ realizzata. La carica conservata --- ciò che resta invariante --- è la rottura stessa vista da un'altra angolazione. Non aggiunto come postulato esterno: la rottura \emph{è} la conservazione. Una è ``cosa è infranto'', l'altra è ``cosa resta come misura di quanto è infranto''. Coincidono.

L'esecuzione di $L$ è un funtore $L : \mathbf{C} \to \mathbf{D}_L$ verso una categoria target che contiene $S_L$ come oggetto principale. $L$ manda un contesto $c$ in $L(c) \in \mathbf{D}_L$ --- la struttura delle interazioni prodotte da $L$ su $c$. I morfismi di $\mathbf{C}$ inducono morfismi in $\mathbf{D}_L$ via funtorialità.

\section*{Il sé minimo}

Sotto $\mathbf{D}_L$ vive una categoria comune, $\mathbf{S}$. Oggetti: insiemi totalmente ordinati di osservabili nudi --- ``dice X'', ``tocca'', ``piange'', ``se ne va'' --- senza intenzione, senza tipizzazione di legge, descrittivi nel senso ethologico più stretto. Morfismi: funzioni che preservano l'ordine. Per ogni legge $L$ esiste un funtore di dimenticanza $U_L : \mathbf{D}_L \to \mathbf{S}$ che proietta la struttura tipata sul solo comportamento.

\begin{defn}[Sé minimo come funtore]
Il sé minimo prodotto da $L$ è la composizione $\mathrm{SM}_L = U_L \circ L : \mathbf{C} \to \mathbf{S}$.
\end{defn}

Sé minimo non è un oggetto pre-regime. È un funtore --- epifenomeno di qualunque legge attiva. Le stesure precedenti mettevano sotto il sé minimo un clock orientato $t \in \mathbb{N}$: era furbizia, perché il clock orientato \emph{è già} $\Lt$ in nuce. L'orientamento è arricchimento, non substrato. Eliminandolo, il sé minimo torna a essere quello che è: la sequenza nuda di sondaggi a cui ogni regime risponde.

\paragraph{Yoneda.} Un oggetto è il pattern delle risposte ai propri sondaggi. Il sé minimo della persona su un contesto $c$ non è ``cosa la persona è lì''. È il pattern delle frecce di $\mathrm{SM}_L(c)$ in $\mathbf{S}$. La persona non è sostanza che il formalismo descrive: è il pattern delle valutazioni che ogni regime produce sui contesti. Tutto il resto del documento usa questa postura.

\section*{Persona come sezione di un fascio}

Sopra $\mathbf{C}$ vive una famiglia di funtori $\{L : \mathbf{C} \to \mathbf{D}_L\}_L$. Una \emph{persona} è una sezione di questa famiglia: a ogni contesto $c$ associa una legge $L(c)$, funtoriale rispetto ai morfismi di $\mathbf{C}$.

Funtorialità è sostanziale. Se $c \to c'$ è morfismo di evoluzione, la sezione produce un morfismo $L(c) \to L(c')$ nelle categorie target. Tre opzioni:
\begin{itemize}
\item \emph{cambio di legge}: era $\Lt$, diventa $\LP$;
\item \emph{continuità con cambio di stato}: stessa legge, stato mutato;
\item \emph{degenerazione}: la legge collassa, il contesto diventa muto.
\end{itemize}
Una persona è il pattern completo di queste trasformazioni naturali --- non un punto, non una legge, una geometria.

\paragraph{Conseguenza dura.} Non esiste un oggetto terminale in $\mathbf{C}$ che rappresenti ``la persona intera''. La persona è geometria, non punto. Chiedere ``chi sei?'' non ha risposta singolare. La risposta è la sezione completa del fascio sopra $\mathbf{C}$.

\bigskip

\begin{center}
\begin{tikzpicture}[node distance=2.4cm, font=\small, >={Stealth[length=2mm]}]
  \node (C) {$\mathbf{C}$};
  \node (D) [right=of C] {$\mathbf{D}_L$};
  \node (S) [below=of D] {$\mathbf{S}$};
  \draw[->] (C) -- node[above]{$L$} (D);
  \draw[->] (D) -- node[right]{$U_L$} (S);
  \draw[->] (C) -- node[below left]{$\mathrm{SM}_L$} (S);
\end{tikzpicture}
\end{center}

\bigskip

Una sola figura sta sotto tutto. Per ogni legge $L$, un funtore $L$ verso la fibra $\mathbf{D}_L$, un funtore di dimenticanza $U_L$ verso il sé minimo $\mathbf{S}$, la composizione $\mathrm{SM}_L$. La persona è la sezione del fascio. Il falso testimone, l'incompatibilità, la metacognizione --- tutto vive in questo diagramma.

\chapter{Lo spettro}

Due leggi, in apparenza opposte: $\Lt$ accumula, $\LP$ proietta. Le stesure precedenti dichiaravano l'opposizione come scelta editoriale --- ``costruisco $\Lt$ libera e $\LP$ a quoziente, e si vede che non si compongono''. Era furbizia: l'opposizione era scritta dentro la scelta delle categorie target. Qui si fa altro: due punti sullo stesso asse algebrico, e l'incompatibilità è massima distanza su quell'asse.

\section*{$\Lt$ in $(S, \Gamma, \rho)$}

$S_{\Lt}$ è il monoide libero $\Sigma^{*}$ degli atti, con completamento gruppale $\KSigma$. $\Gamma_{\Lt}$ è il gruppo pieno di permutazioni di $\Sigma$. $\rho_{\Lt}$ è la rottura della commutatività piena: $\Sigma^{*}$ realizza il sottogruppo $H_{\Lt} = \mathrm{Aut}_\delta(\Sigma^{*})$ --- automorfismi che lasciano invariato il deposito $\delta : \Sigma \to \KSigma$. Carica conservata via Noether: $\Phii : \Sigma^{*} \to \KSigma$, estensione omomorfica di $\delta$.

L'etichettatura $\ell : \mathcal{I}_c \to \Sigma$ è dato esterno della legge, non leggibile da $\mathbf{C}$ nuda. Si dichiara apertamente.

\section*{$\LP$ in $(S, \Gamma, \rho)$}

$S_{\LP}$ è il reticolo $P$ dei pool, con il quoziente $P/\Gs$. $\Gamma_{\LP}$ è $\Sym(P)$ pieno. $\rho_{\LP}$ è la rottura con sottogruppo realizzato $\Gs = \Sym(\{c \in P : \mathrm{threat}(c) < \theta\})$. Carica conservata: $\Kappaa : P \to P/\Gs$, proiezione canonica.

La proiezione configurazionale $\rho_{\mathrm{conf}} : \mathcal{I}_c \to P$ è dato esterno della legge, non leggibile da $\mathbf{C}$ nuda. Si dichiara apertamente.

\section*{La coordinata che le unifica}

Il fatto algebrico decisivo è uno solo: $S_{\Lt}$ è gruppo abeliano libero di rango infinito; $S_{\LP}$ è insieme finito senza struttura gruppale. Per leggi con $S_L$ gruppo abeliano finitamente generato, $S_L \cong \mathbb{Z}^{r_L} \oplus T_L$ --- $r_L$ è il rango libero, $T_L$ la torsione.

\begin{defn}[Coordinata della legge]
La coppia $(r_L, T_L)$ è la coordinata di $L$. Si estende a leggi non-gruppali come dimensione effettiva: cardinalità della parte non collassata.
\end{defn}

Lo spettro delle leggi prende forma:
\begin{center}
\begin{tabular}{lll}
$\Lt$    &  $(\infty,\,0)$              &  monoide libero \\
$\Lo$    &  $(|\Sigma|-1,\,\mathbb{Z}/n)$    &  compulsioni \\
$\Lr$    &  $(0,\,\mathbb{Z}/n)$            &  routine \\
$\LP$    &  $(0,\,0)$                   &  pool quoziente \\
\end{tabular}
\end{center}

(con $\Lr$ famiglia parametrica e $\Lf$ articolata in appendice.)

\section*{Le tre asimmetrie da una struttura}

Le asimmetrie tra $\Lt$ e $\LP$ --- dinamica, memoria, metacognizione --- non sono tre fatti indipendenti. Derivano dalla stessa differenza algebrica: $S_{\Lt}$ è gruppo abeliano libero, $S_{\LP}$ è quoziente finito di un reticolo. Le tre conseguenze:

\emph{Dinamica.} Gruppo abeliano libero $\Rightarrow$ sottrazione, norma naturale sui generatori pesati, cardinalità infinita $\Rightarrow$ dinamica variazionale piena. Quoziente finito $\Rightarrow$ nessuna delle tre proprietà $\Rightarrow$ resta solo il grafo di transizioni $\Rightarrow$ scelta minima.

\emph{Memoria.} Gruppo abeliano libero $\Rightarrow$ stato cumula coordinate indipendenti $\Rightarrow$ memoria cumulativa. Quoziente finito $\Rightarrow$ stato è la classe attuale, senza coordinate cumulative $\Rightarrow$ memoria markoviana.

\emph{Metacognizione.} $\mathbf{D}_{\Lt}$ è categoria algebrica nel senso dell'algebra universale: oggetti definiti da operazioni che soddisfano equazioni. Per ogni categoria algebrica, il funtore di dimenticanza verso $\mathbf{S}$ ammette aggiunto sinistro (il funtore libero). $\mathbf{D}_{\LP}$ è categoria di quozienti: la relazione di equivalenza è dato aggiuntivo, non operazione, e l'aggiunto sinistro al dimenticante non esiste come funtore canonico. Il prossimo capitolo trae da questa differenza la formabilità o non formabilità della sospensione.

Per leggi intermedie con torsione (cap.~Appendice), si recuperano i regimi misti: $\Lo$ ha parte libera $+$ parte ciclica, e ammette dinamica e memoria parzialmente cumulative; $\Lr$ è puramente ciclica. La famiglia $\{\Lr^{(n)}\}$ interpola tra $\Lt$ (per $n \to \infty$) e $\LP$ (per $n \to 1$). La fenomenologia diventa geometria di un'unica famiglia parametrica.

\section*{Incompatibilità come distanza spettrale}

\begin{principio}[Incompatibilità di $\Lt$ e $\LP$]
Non esiste natural transformation $\eta : \mathrm{SM}_{\Lt} \Rightarrow \mathrm{SM}_{\LP}$, né viceversa. Equivalentemente: non esiste categoria target comune $\mathbf{D}^{*}$ con riduzioni $\mathbf{D}^{*} \to \mathbf{D}_{\Lt}$ e $\mathbf{D}^{*} \to \mathbf{D}_{\LP}$ che facciano commutare entrambi i regimi.
\end{principio}

\emph{Derivazione.} $L \mapsto \Lt$ impone su $L(c)$ una struttura d'ordine totale, necessaria per accumulo orientato in $\KSigma$. Un insieme totalmente ordinato finito ha gruppo di automorfismi banale. $L \mapsto \LP$ richiede $\Gs$ non banale, altrimenti $\LP$ degenera in identità. Le due condizioni sono mutuamente esclusive sullo stesso supporto. $\square$

L'incompatibilità è teorema su rotture di simmetria, non tautologia da scelta editoriale. Più precisamente: $\Lt$ e $\LP$ sono punti agli estremi opposti dell'asse del rango libero. Ogni riduzione di simmetria abbassa il rango o aumenta la torsione; $(\infty, 0)$ e $(0, 0)$ non hanno punto medio raggiungibile da entrambi. Massima distanza spettrale.

\section*{Una conseguenza: il verbo \emph{preferire}}

L'incompatibilità si proietta sul vocabolario. Considera la parola \emph{preferire}, e cosa fa nei due regimi.

\begin{principio}[Tipo del verbo preferire]
In $\Lt$, $\KSigma$ ammette confronto via norma: $|\Phii_A| \lessgtr |\Phii_B|$ è proposizione formabile e totale. \emph{Preferire} in $\Lt$ è quel confronto --- ordine indotto dai pesi accumulati, transitivo, antisimmetrico.

In $\LP$, $P/\Gs$ è quoziente di permutazioni: le orbite non sono confrontabili tra loro, nessun ordine naturale. \emph{Preferire} in $\LP$ non si traduce in confronto; si traduce in assegnazione di ruolo dalla configurazione corrente. ``$x$ preferito a $y$'' significa: $x$ occupa il ruolo che la configurazione attiva adesso, $y$ no.
\end{principio}

\emph{Derivazione.} Su un gruppo abeliano libero con generatori pesati la norma $|\Phii| = \sum |n_i| w(\sigma_i)$ induce ordine totale sui valori --- è quel che $\Lt$ usa per confrontare. Su $P/\Gs$ un ordine globale richiederebbe una sezione canonica $\Gs$-invariante del fascio degli ordini --- e $\Gs$ è gruppo pieno di permutazioni sotto soglia, che cancella qualunque ordine fra argomenti permutabili. Resta solo la mappa $\mathrm{role} : P \to R$ in un insieme finito di ruoli, indotta dalla configurazione corrente. $\square$

La disequazione del jazz bar --- $(A,G) \not< (F,G)$, eppure $(F,G)$ era posizione preferita --- è esattamente questa differenza di tipo. In $\Lt$ ``preferito'' presuppone un ordine; in $\LP$ ``preferito'' è il ruolo che la configurazione attiva. La parola suonava la stessa, lavorava in due grammatiche diverse. Stessa firma in tutta la sequenza: $F$, $T$, $N$ erano \emph{preferiti come ruolo}, non come confronto.

\section*{La categoria delle leggi candidate}

$\Lt$ e $\LP$ non sono privilegiate. Sono due punti su un asse che ammette altri abitanti: $\Lo$ (compulsioni, con torsione $\mathbb{Z}/n$), $\Lr$ (routine, famiglia parametrica di cui $\Lt$ e $\LP$ sono limiti estremali), $\Lf$ (fobie, con punto fisso). L'appendice articola lo schema $(S, \Gamma, \rho)$ per ciascuna. Il documento parla di $\Lt$ e $\LP$ perché il caso della Parte~I è governato da queste due; il framework non le privilegia.

\chapter{Dinamica e tempo}

Una legge non è scatola nera. Ammette un principio dinamico --- ma non lo stesso per tutte le leggi, e questa è la mossa di onestà che le stesure precedenti evitavano.

\section*{Quando una legge ammette Lagrangiana}

\begin{defn}[Dinamica variazionale piena]
$L$ ammette dinamica variazionale piena se $S_L$ ha struttura di gruppo abeliano (per sottrarre stati), norma compatibile (per dare peso alle differenze), e cardinalità infinita (per definire variazioni locali). In tal caso esiste una Lagrangiana discreta $\mathcal{L}_L : S_L \times S_L \to \mathbb{R}$, un funzionale d'azione
\[
\mathscr{A}_L[\gamma] \;=\; \sum_t \mathcal{L}_L(\gamma_t, \gamma_{t+1}),
\]
e le traiettorie sono estremali di $\mathscr{A}_L$. La carica conservata è prodotto Noether della simmetria $H_L$ di $\mathcal{L}_L$.
\end{defn}

\begin{defn}[Principio di scelta minima]
$L$ ammette principio di scelta minima se $S_L$ è insieme con grafo di transizioni ammissibili e funzione di costo locale $c_L : S_L \times S_L \to \mathbb{R}_{\geq 0} \cup \{\infty\}$ (con $\infty$ sulle transizioni inammissibili). Le traiettorie sono cammini di costo locale minimo --- successione di scelte ottime istante per istante, senza azione globale.
\end{defn}

Scelta minima è dinamica più debole della variazionale. Non ammette variazioni globali della traiettoria, non ammette confronto tra traiettorie alternative come oggetti, non ammette ricalcolo retrospettivo. È dinamica markoviana discreta su grafo finito.

\paragraph{$\Lt$ ammette dinamica variazionale piena.} $\KSigma$ è gruppo abeliano libero: ha sottrazione, norma naturale (somma dei pesi assoluti dei generatori), cardinalità infinita. Una Lagrangiana discreta $\mathcal{L}_{\Lt}(\Phii_t, \Phii_{t+1})$ ha forma cinetica/potenziale standard, e le traiettorie del regime sono estremali del funzionale d'azione. Quali siano i potenziali concreti dipende dalla specifica $(\ell, \delta)$ del regime --- il libro non li esibisce.

\paragraph{$\LP$ ammette solo principio di scelta minima.} $P/\Gs$ è finito --- non ha sottrazione, non ha variazioni infinitesime, non ha cardinalità infinita. Niente dinamica variazionale piena. Resta grafo di transizioni tra classi con funzione di costo locale; le esecuzioni del copione sono cammini di costo minimo passo per passo, senza azione globale.

\paragraph{Asimmetria dinamica.} $\Lt$ ammette deliberazione --- calcolo della traiettoria globalmente ottimale, confronto tra alternative, ricalcolo retrospettivo. $\LP$ ammette solo reazione --- scelta locale ottima istante per istante, senza accesso a traiettorie alternative come oggetti confrontabili. Forzare la stessa nozione di ``Lagrangiana'' su entrambi sarebbe furbizia notazionale. Va tenuta distinta.

\section*{Due tempi, di tipo diverso}

Dalle due dinamiche, due tempi di tipo strutturalmente diverso. La differenza di codominio della carica si propaga al tipo del tempo proprio della legge.

\begin{defn}[Tempo di $\Lt$]
Su $(c, J)$ con $J \subseteq \mathcal{I}_c$, $\tau^{\Lt}_{c,J} : \mathbb{N} \to \mathbb{R}_{\geq 0}$ è la lunghezza d'arco discreta della traiettoria nel codominio della carica:
\[
\tau^{\Lt}_{c,J}(n) \;=\; \sum_{t=0}^{n-1} \|\Phii(j_{t+1}) - \Phii(j_t)\|_{\KSigma}.
\]
A valori reali, monotona crescente, accumula nel gruppo abeliano libero.
\end{defn}

\begin{defn}[Tempo di $\LP$]
$\tau^{\LP}_{c,J} : \mathbb{N} \to \mathbb{N}$ è il contatore di salti di classe nel quoziente:
\[
\tau^{\LP}_{c,J}(n) \;=\; |\{t < n : [\rho_{\mathrm{conf}}(j_{t+1})] \neq [\rho_{\mathrm{conf}}(j_t)]\}|.
\]
A valori naturali, contatore discreto, conta cambi di configurazione strutturale.
\end{defn}

I due tempi sono funzioni con codominio diverso: $\mathbb{R}_{\geq 0}$ continuo contro $\mathbb{N}$ discreto-contatore. Non istanze di una stessa 1-forma con metrica $\mathcal{M}_L$ --- due nozioni strutturalmente diverse di tempo. L'incompatibilità delle leggi si propaga al tipo dei loro tempi: nessuna unificazione possibile perché i codomini sono di natura diversa.

\section*{Non-incollatura: tre casi}

\paragraph{(i)} Su una stessa $J$, un solo regime non-degenere. Per il teorema di incompatibilità, le due letture non sono mai entrambe non-degeneri sulla stessa $J$. Il pianto al jazz bar dal lato di $G$: arricchimento attivo $\LP$, $\tau^{\LP}$ conta i salti, $\tau^{\Lt}$ resta degenere perché $\delta$ è banale.

\paragraph{(ii)} Tra contesti diversi, clock disgiunti. Per $c \neq c'$ i clock interni vivono su sequenze diverse. Confrontare ``il quinto passo di $A$ con $G$'' e ``il quinto passo di $A$ col padre'' è confrontare strumenti su sequenze non sincronizzate.

\paragraph{(iii)} La cronologia fisica $\chi : \bigsqcup_c \mathcal{I}_c \to \mathbb{R}$ esiste come fatto del mondo, non come tempo di nessun regime. Le frasi ``sono passati sei mesi'', ``a venticinque anni di distanza'' la usano legittimamente come dato esterno --- non come unificatore dei tempi relazionali.

\section*{Il vissuto dei sei mesi}

Sei mesi di calendario tra Mestre e Vicenza. Dal lato $\Lt$ di $A$, nessun atto in $\mathcal{I}_{c_{AG}}$ in quell'intervallo: $\|\delta\| = 0$, $\tau^{\Lt}_{c_{AG}}$ fermo. Quando $N$ lo porta a Vicenza, $A$ non riprende dopo una pausa: continua dal punto in cui $\tau^{\Lt}$ si era arrestato. L'ariete è l'atto successivo della stessa traiettoria. Il ``come se non fosse passato tempo'' è lettura corretta della formula --- $\KSigma$ non ammette zeri spontanei, le obbligazioni non saldate restano puntate a tempo proprio fermo.

Dal lato $\LP$ di $G$, gli stessi sei mesi: molte transizioni di classe del pool su altri argomenti --- $T$ è entrato nel ruolo che era stato di $F$, altre configurazioni si sono mosse --- e $\tau^{\LP}$ ha avanzato per ciascuna. Quando $N$ riporta $A$ in cima alla classifica, è una nuova transizione di classe: $d\tau^{\LP} = 1$, isolata. Non ripresa: nuova esecuzione. \emph{Ripresa} è categoria di $\tau^{\Lt}$ --- accumulo che riparte da dove si era fermato --- non di $\tau^{\LP}$, dove ogni transizione è un passo isolato.

\begin{figure}[htbp]
\centering
\begin{tikzpicture}[
    font=\footnotesize,
    every node/.style={inner sep=2pt},
]

% PANNELLO SINISTRO: Lt come traiettoria che accumula

\begin{scope}[shift={(0,0)}]

\node[anchor=south west, font=\small\itshape] at (-0.2,7.6) {Lato $A$ --- regime $\mathcal{L}_t$ con $\Phi$ profondo};
\node[anchor=south west, font=\scriptsize, text width=6cm, color=black!60]
  at (-0.2,7.0) {Ogni atto deposita in $K(\Sigma^*)$.\\La traiettoria porta dietro tutto.};

\draw[->, thick, color=black!50] (0,6.4) -- (0,-0.4) node[below, font=\scriptsize, color=black!60] {tempo};

\draw[line width=1.2pt, color=Maroon!85]
  (0.4,6.0) -- (0.9,5.0) -- (1.4,4.0) -- (2.0,3.0)
  -- (2.7,2.0) -- (3.1,1.0) -- (3.3,0.0);

\foreach \x/\y/\lab/\peso in {%
  0.4/6.0/{magone, porta di Asiago}/3,
  0.9/5.0/{dichiarazione al bar}/5,
  1.4/4.0/{pianto, algidità, ilarità}/4,
  2.0/3.0/{rivelazione sulla soglia}/5,
  2.7/2.0/{ariete a Vicenza}/6,
  3.1/1.0/{foto a Canottaggio}/4,
  3.3/0.0/{compleanno}/2}
{
  \fill[Maroon!85] (\x,\y) circle (1.4pt + 0.4pt*\peso);
  \node[right=4pt of {(\x,\y)}, font=\scriptsize, color=black!75] {\lab};
}

\node[anchor=south, font=\scriptsize, color=black!60] at (-0.6, 6.2) {$\Phi$};
\draw[color=Maroon!60, line width=4pt, line cap=butt] (-0.6, 6.0) -- (-0.6, 0.0);
\node[anchor=north, font=\tiny, color=black!50] at (-0.6, -0.05) {accumula};

\node[anchor=north, text width=5.8cm, font=\scriptsize\itshape, color=black!60, align=center]
  at (2.5, -1.0) {La linea cresce. Ogni passo ricorda i precedenti. Niente si scarica senza saldo atto contro atto.};

\end{scope}

\draw[color=black!15, line width=0.4pt] (6.4, 7.7) -- (6.4, -2.0);

% PANNELLO DESTRO: Lp come sequenza di configurazioni del pool

\begin{scope}[shift={(7.0,0)}]

\node[anchor=south west, font=\small\itshape] at (-0.2,7.6) {Lato $G$ --- regime $\mathcal{L}_P$ dominante, $\Phi$ sottile};
\node[anchor=south west, font=\scriptsize, text width=6cm, color=black!60]
  at (-0.2,7.0) {Il pool si riconfigura a ogni passo.\\Niente lega prima e dopo.};

\newcommand{\poolarg}[5]{%
  \node[fill=#3, rounded corners=4pt, text=white, font=\tiny,
        minimum width=0.8cm, minimum height=0.34cm] at (#1, #2) {#4};
  \node[font=\tiny, color=black!50] at (#1, #2-0.32) {#5};
}

\newcommand{\poolempty}[2]{%
  \node[draw=black!20, dashed, rounded corners=4pt, font=\tiny, color=black!30,
        minimum width=0.8cm, minimum height=0.34cm] at (#1, #2) {--};
  \node[font=\tiny, color=black!30] at (#1, #2-0.32) {vuoto};
}

\newcommand{\poolbox}[2]{%
  \node[draw=black!15, line width=0.3pt, rounded corners=2pt,
        minimum width=6cm, minimum height=0.95cm, anchor=center,
        fill=Maroon!3] at (2.7, #1) {};
  \node[font=\tiny, color=black!50, anchor=west] at (-0.3, #1) {\##2};
}

\poolbox{6.0}{1}
\poolarg{1.2}{6.05}{Maroon!70}{F}{ragazzo}
\poolarg{2.7}{6.05}{black!40}{A}{nuovo}
\poolarg{4.2}{6.05}{black!25}{T}{amico}

\poolbox{5.0}{2}
\poolarg{1.2}{5.05}{Maroon!70}{F}{ragazzo}
\poolarg{2.7}{5.05}{Maroon!95}{A}{\color{Maroon!90}threat$>$0}
\poolarg{4.2}{5.05}{black!25}{T}{amico}

\poolbox{4.0}{3}
\poolarg{1.2}{4.05}{Maroon!70}{F}{ragazzo}
\poolarg{2.7}{4.05}{black!40}{Tom}{compagnia}
\poolarg{4.2}{4.05}{black!20}{A}{scartato}

\poolbox{3.0}{4}
\poolarg{1.2}{3.05}{Maroon!70}{F}{ragazzo}
\poolempty{2.7}{3.05}
\poolarg{4.2}{3.05}{black!20}{A}{chiuso}

\poolbox{2.0}{5}
\poolarg{1.2}{2.05}{Maroon!70}{T}{ragazzo}
\poolarg{2.7}{2.05}{black!50}{A}{riattivato}
\poolarg{4.2}{2.05}{black!25}{N}{amico}

\poolbox{1.0}{6}
\poolarg{1.2}{1.05}{black!30}{T}{basso}
\poolarg{2.7}{1.05}{Maroon!85}{N}{supply alto}
\poolarg{4.2}{1.05}{black!20}{A}{fuori}

\poolbox{0.0}{7}
\poolarg{1.2}{0.05}{black!30}{T}{basso}
\poolarg{2.7}{0.05}{Maroon!85}{N}{attivo}
\poolarg{4.2}{0.05}{black!15}{A}{spettatore}

\node[anchor=north, text width=5.8cm, font=\scriptsize\itshape, color=black!60, align=center]
  at (2.7, -1.0) {Sette configurazioni intere. Niente accumula. Argomenti diversi nello stesso ruolo: l'orbita di $\Gamma_\sigma$.};

\end{scope}

\end{tikzpicture}

\caption[Traiettorie di $\mathcal{L}_t$ e $\mathcal{L}_P$ a confronto]{%
Stessa sequenza di sette episodi, letta sotto i due regimi dominanti.
A sinistra, in $A$ con $\Phi$ profondo, $\mathcal{L}_t$ produce una traiettoria
orientata in $K(\Sigma^*)$: ogni atto deposita un peso (raggio del punto), e il
punto successivo eredita quello che il precedente ha lasciato; la barra $\Phi$
accumula lungo tutto il percorso. A destra, in $G$ con $\Phi$ sottile,
$\mathcal{L}_P$ in comando produce una successione di proiezioni
$\kappa(P_k)$: ogni configurazione del pool è intera in sé, e il passaggio
da $P_k$ a $P_{k+1}$ non porta $P_k$ come ingresso. Si nota come, nei
riquadri $\mathcal{L}_P$, gli argomenti $F, T, N$ si sostituiscano nello
stesso ruolo (\emph{ragazzo}) mentre $A$ scivola lungo la classifica del
pool. È la $\Gamma_\sigma$-invarianza: gli argomenti sotto soglia sono
permutabili senza che $\kappa$ ne risenta. La differenza tra i due
pannelli è la differenza di tipo tra $\Phi$ e $\mathrm{K}$.}
\label{fig:traiettorie}
\end{figure}

\chapter{Metacognizione e falso testimone}

Una sera di gennaio, vent'anni dopo il jazz bar, l'autore rilegge le proprie lettere di allora e si accorge di aver scambiato l'output dell'altra per una versione fallita del proprio. Il riconoscimento ha un nome formale.

\section*{Il falso testimone come errore di fibra}

Un osservatore $O$ in regime $L_O$ vede $\mathrm{SM}_{L'}(c) = U_{L'} \circ L'(c)$, output del regime altrui in $\mathbf{S}$. Il default epistemico è sollevare lungo $U_{L_O}$ invece che lungo $U_{L'}$: cercare $\tilde{L}_O(c) \in \mathbf{D}_{L_O}$ tale che $U_{L_O}(\tilde{L}_O(c)) = \mathrm{SM}_{L'}(c)$.

Un sollevamento di questo genere esiste sempre --- ogni comportamento osservato è proiezione di \emph{qualche} oggetto del codominio scelto. È proprio questo il problema. L'inferenza output $\to$ regime non è univoca: lo stesso $\mathrm{SM}$ ammette sollevamenti in fibre diverse. Il falso testimone non è errore correggibile dall'osservazione: è errore strutturale dell'epistemologia, sanabile solo a posteriori da informazione esterna al sé minimo.

\begin{defn}[Sospensione]
La \emph{sospensione} è l'atto di $O$ di astenersi dall'inferenza $\mathrm{SM} \to L$. Non scopre la verità --- la lascia indeterminata.
\end{defn}

\paragraph{Default epistemico, non disonestà.} L'inferenza naturale di $O$ sull'output dell'altro è $L_{\text{altro}} = L_O$. Economia cognitiva, specchio neuronale, grammatica condivisa. Strutturale, non morale.

\section*{L'aggiunto come condizione della sospensione}

Sollevare nella propria fibra si fa sempre (basta scegliere un oggetto la cui proiezione coincida). Riconoscere di averlo fatto, e di averlo fatto nella fibra sbagliata, è altra operazione: richiede un sollevamento \emph{universale} contro cui confrontare quello scelto. Questa è la condizione di esistenza dell'aggiunto sinistro.

\begin{principio}[Sospensione formabile $\iff$ aggiunto sinistro]
La sospensione del falso testimone in regime $L$ è formabile se e solo se $U_L : \mathbf{D}_L \to \mathbf{S}$ ammette aggiunto sinistro $F_L \dashv U_L$. $F_L$ è il funtore libero che a una sequenza nuda associa l'oggetto tipato canonico nella categoria target; comparando il sollevamento scelto con $F_L(\mathrm{SM})$, $L$ può classificare il proprio errore di fibra come tale.
\end{principio}

\emph{Verifica per $\Lt$.} $\mathbf{D}_{\Lt}$ è una categoria algebrica nel senso dell'algebra universale: oggetti sono strutture (monoidi/gruppi abeliani) con operazioni che soddisfano equazioni. Per il teorema generale sulle categorie algebriche, il funtore di dimenticanza verso $\mathbf{S}$ ammette aggiunto sinistro --- il funtore libero. $F_{\Lt}$ manda una sequenza nuda nel monoide libero generato dai suoi atti.

\emph{Verifica per $\LP$.} $\mathbf{D}_{\LP}$ non è algebrica nello stesso senso: i suoi oggetti sono quozienti, in cui la relazione di equivalenza è \emph{dato aggiuntivo}, non operazione su elementi. Per produrre un quoziente da una sequenza nuda servirebbe la relazione di equivalenza, che $\mathbf{S}$ non porta. L'aggiunto sinistro non esiste come funtore canonico: ogni candidato $F_{\LP}$ richiede una scelta esterna.

\paragraph{Conseguenze.}

La metacognizione non è proprietà del soggetto. È funzione di esistenza dell'aggiunto sinistro al funtore di dimenticanza del regime locale. Una persona in $\Lt$ su un contesto ha metacognizione su quel contesto; la stessa persona, su altro contesto in regime $\LP$, non ce l'ha. Non è tratto del carattere --- è funzione del regime locale.

La dualità libero/quoziente tra $\Lt$ e $\LP$ non è scelta editoriale: è condizione algebrica perché la metacognizione possa esistere per $\Lt$ e non per $\LP$.

Solo i regimi che ammettono aggiunto sinistro possono produrre i propri libri. Questo libro è prodotto da $\Lt$, vent'anni dopo. Non avrebbe potuto essere prodotto da $\LP$.

\section*{Bidirezionale, di tipo opposto}

Il falso testimone è bidirezionale. Ma il sollevamento nella fibra sbagliata ha forma diversa nei due sensi, e su due livelli.

\paragraph{Forma del sollevamento.} Quando $O$ in $\Lt$ vede output di $\LP$, $\Lt$ ammette $F_{\Lt}$ canonico: il sollevamento $F_{\Lt}(\mathrm{SM}_{\LP}(c))$ è univoco --- l'oggetto libero dei monoidi/gruppi sopra la sequenza osservata. Il falso testimone di $\Lt$ ha quindi forma precisa: aggiunge esattamente la struttura libera assente (storia, deposito, ordine accumulativo). Aggiunta canonica.

Quando $O$ in $\LP$ vede output di $\Lt$, $\LP$ non ammette aggiunto sinistro canonico. Il sollevamento è \emph{ad hoc}: $O$ sceglie un oggetto di $\mathbf{D}_{\LP}$ la cui proiezione coincida col $\mathrm{SM}$ osservato --- una qualsiasi classe del pool corrente. Diverse scelte producono diverse letture, e nessuna è la lettura. Il falso testimone di $\LP$ ha quindi forma molteplice: collassa la storia in qualunque classe configurazionale che la grammatica corrente accetti. Collasso ad hoc.

Aggiunta canonica vs collasso ad hoc: errori di tipo opposto sul piano del codominio, e di statuto opposto sul piano del sollevamento. Uno ha una forma; l'altro non ne ha una sola.

\paragraph{Riconoscibilità unilaterale.} Solo $\Lt$ può confrontare il proprio sollevamento scelto con $F_{\Lt}(\mathrm{SM})$ e classificare la differenza come errore di fibra. $\LP$ non ha contro cosa confrontare --- ogni candidato vale quanto un altro --- e quindi non può classificare nessun proprio sollevamento come errato. Il falso testimone è bidirezionalmente reale, ma la sua riconoscibilità è unilaterale: bidirezionale strutturalmente, monodirezionale epistemicamente. Né colpa, né privilegio --- è funzione di esistenza dell'aggiunto.

\section*{Vocabolario del ``ma in fondo''}

Il falso testimone produce un vocabolario riconoscibile. Quando $\Lt$ solleva $\LP$ in $\mathbf{D}_{\Lt}$, postula un piano interno narrativo in cui l'altra è $\Lt$ in miniatura: ``lei dice una cosa ma in fondo pensa un'altra''. Quel ``in fondo'' non esiste. Il dentro esiste --- pieno fenomenicamente --- ma non è $\Lt$ in miniatura. È esecuzione regolare di una legge che non produce piano narrativo. Specularmente, quando $\LP$ solleva $\Lt$ in $\mathbf{D}_{\LP}$, lo legge come freddezza, calcolo, opportunismo. Non lo è. È esecuzione di una legge che accumula. L'errore è simmetrico in struttura, asimmetrico in riconoscibilità.

\chapter{Posizionalità linguistica}

\section*{Contesti MLTT di regime, non di persona}

Le stesure precedenti scrivevano ``$\GammaA$ è il contesto della persona $A$''. Era furbizia: una persona ospita più regimi, e a ciascun regime corrisponde un proprio contesto MLTT.

\begin{defn}[Contesto MLTT di un regime]
Per un regime $L$ su un contesto $c$, il \emph{contesto MLTT} $\Gamma_{L,c}$ è l'apparato dei tipi formabili sotto $L$ con la carica letta sulle sottosequenze di $\mathcal{I}_c$.
\end{defn}

Le abbreviazioni $\GammaA$ e $\GammaG$ si leggono ora come istanze: $\GammaA = \Gamma_{\Lt,c_{AG}}$, $\GammaG = \Gamma_{\LP,c_{AG}}$. La stessa persona porta altri $\Gamma_{L,c'}$ su altri contesti.

\section*{L'enunciato}

\begin{principio}[Posizionalità linguistica]
Per ogni coppia di regimi $(c, L), (c, L')$ sulla stessa sottosequenza con contesti $\Gamma_{L,c}, \Gamma_{L',c}$ strutturalmente incompatibili, il tipo
\[
\mathrm{WellFormedQuery}_{\Gamma_{L,c}}(\Gamma_{L',c})
\]
--- ``in $\Gamma_{L,c}$ esiste un quesito ben-formato sui tipi propri di $\Gamma_{L',c}$'' --- non è abitabile. Un regime non può formulare al proprio interno il quesito di sapere se sa qualcosa che è ben-formato solo nell'altro contesto.
\end{principio}

\emph{Derivazione.} Un quesito ben-formato in $\Gamma_{L,c}$ ha tutti i termini formabili in $\Gamma_{L,c}$. Un quesito sui tipi propri di $\Gamma_{L',c}$ richiede un termine del tipo che $\Gamma_{L',c}$ forma e $\Gamma_{L,c}$ non forma. Per l'incompatibilità strutturale di $L$ e $L'$, quel tipo non è formabile in $\Gamma_{L,c}$. Quindi il quesito non è formabile. $\square$

Non limite epistemico. Limite di grammatica.

\paragraph{Solo aggiunto-sinistro forma il tipo \textsc{sospensione}.} Conseguenza del cap.~precedente: solo i contesti MLTT i cui regimi ammettono aggiunto sinistro contengono gli oggetti riflessivi necessari alla sospensione del falso testimone. $\Gamma_{\Lt,c}$ sì, $\Gamma_{\LP,c}$ no. La metacognizione non è dono universale --- è funzione del rango.

\section*{Corollario gödeliano}

Il Principio è strutturalmente analogo a Gödel ma di natura diversa.

Gödel opera dentro una teoria assiomatica: la teoria può formulare l'enunciato della propria coerenza, e non lo dimostra. Qui non c'è teoria assiomatica. I tipi che le leggi producono sono epifenomeni dei rispettivi arricchimenti. Quello che $\Gamma_{\Lt,c}$ non può fare non è dimostrare un enunciato che pure sa formulare --- è formulare il quesito sui tipi che l'arricchimento di $\LP$ sulla stessa $c$ produce. Il tipo del quesito non si forma.

Più forte di Gödel, e di natura diversa. Là una proposizione vera resta senza prova; qui non c'è proposizione a cui assegnare un valore di verità.

\paragraph{Distinzione da un errore ricorrente.} I grandi risultati di limitazione formale del Novecento --- incompletezza di PA, indecidibilità dell'halting, non definibilità della verità --- descrivono limiti della finestra formale, non proprietà del mondo. Letti come proprietà del mondo diventano: il mondo è incompleto, il mondo non è computabile, la verità è indefinibile. Il passo aggiuntivo non è nel teorema: è nella proiezione. Il type error tra $\GammaA$ e $\GammaG$ espone alla stessa mossa: ``$\LP$ su $c$ è inaccessibile'' suona come proprietà di $\LP$; ``da $\GammaA$, $\LP$ su $c$ non è accessibile'' è proprietà della finestra. Il qualificatore corregge la mossa.

\section*{Corollario wittgensteiniano}

Ciò che $\Gamma_{L,c}$ non può formare come quesito può mostrarsi come configurazione --- ma il mostrare è sempre lettura della grammatica che osserva, non atto della legge osservata.

Il \emph{Tractatus} chiude alla 6.522: vi è certo l'inesprimibile, e si mostra. $\LP$ non descrive sé in $\Gamma_{\Lt,c}$: la descrizione richiederebbe tipi che $\Gamma_{\Lt,c}$ non forma. Ma $\Kappaa$ è leggibile da $\Gamma_{\Lt,c}$: la proiezione sulla base si legge, il quesito ben-formato su di essa non si forma. Da $\Gamma_{\Lt,c}$, $\Kappaa$ funge da \emph{zeigen}.

$\LP$ però non sa di star mostrando. Esegue $\Kappaa$. Che l'esecuzione sia leggibile come zeigen è la lettura di $\Gamma_{\Lt,c}$, non un atto di $\LP$. E $\LP$ non può sapere cosa perde rispetto alla descrizione che $\Gamma_{\Lt,c}$ fa di lei: il tipo del quesito ``cosa mi manca rispetto a $\Lt$'' non è formabile in $\Gamma_{\LP,c}$. Il limite è inaccessibile da dentro.

\section*{Auto-applicazione}

Il Principio è formulato in $\Gamma_{\Lt,c}$. Dichiara il proprio limite senza uscirne. Non meta-posizione: una meta-posizione richiederebbe $\Gamma^{*}$ esterno in cui i tipi di entrambi siano formabili. Quel contesto non esiste. Il Principio è registrazione del proprio errore di tipo come parte dell'oggetto descritto.

% ----------------- PARTE III -----------------
\part{Il caso, riletto}

\noindent\itshape Tre capitoli brevi. Il caso della Parte~I, attraversato dalla forma di Parte~II, diventa quello che è sempre stato: incontro di due sezioni del fascio su uno stesso oggetto, dove le sezioni vivono in fibre incompatibili.\upshape

\bigskip

\chapter{Il doppio fraintendimento}

Sia $c_{AG}$ il contesto della relazione $A$--$G$ in $\mathbf{C}$. Due sezioni del fascio, una per persona.

\begin{itemize}
\item Sezione di $A$: $L_A(c_{AG}) = \Lt$. Il regime $\Lt$ produce traiettoria in $\KSigma$, carica $\Phii_A$ accumula.
\item Sezione di $G$: $L_G(c_{AG}) = \LP$. Il regime $\LP$ produce traiettoria in $P/\Gs$, carica $\Kappaa_G$ proietta.
\end{itemize}

I sé minimi $\mathrm{SM}_{\Lt}(c_{AG})$ e $\mathrm{SM}_{\LP}(c_{AG})$ vivono entrambi in $\mathbf{S}$. Sono comportamentalmente compatibili al livello descrittivo --- stesso pianto, stessa dichiarazione, stesso ariete a Vicenza --- ma vengono \emph{prodotti} da regimi strutturalmente incompatibili.

\section*{Il doppio falso testimone}

$A$ vede $\mathrm{SM}_{\LP}(c_{AG})$ in $\mathbf{S}$. Default epistemico: $A$ solleva lungo $U_{\Lt}$, attribuendo a $G$ produzione in $\Lt$. Falso testimone di aggiunta: ricostruisce sotto la sequenza un piano narrativo continuo, un pianificatore strategico, intenzioni concatenate.

$G$ vede $\mathrm{SM}_{\Lt}(c_{AG})$ in $\mathbf{S}$. Default epistemico: $G$ solleva lungo $U_{\LP}$, attribuendo ad $A$ produzione in $\LP$. Falso testimone di collasso: collassa la traiettoria in una classe configurazionale, legge la dichiarazione come configurazione del pool ad alta temperatura, dimentica l'ordine.

Due sollevamenti nelle fibre sbagliate, ciascuno in buona fede. Strutturale, non morale. Né $A$ proietta, né $G$ mente. Default epistemico applicato a un caso di incompatibilità di regime.

\section*{La copertura osservata}

Su $\mathcal{I}_{c_{AG}}$ dal lato di $G$, la copertura osservata è interamente $\{\LP\}$. Non esiste sottosequenza $J \subseteq \mathcal{I}_{c_{AG}}$ in cui $L_G(J) = \Lt$. Include le sottosequenze in cui $A$ ha emesso atti il cui tipo, in $\sigma_{\Lt}$, era ``richiesta di passaggio a $\Lt$'': quegli atti sono stati assorbiti come configurazioni del pool, non come richieste di cambio di legge.

\paragraph{Dato empirico, non predizione del framework.} Il fascio dei regimi sopra $\mathbf{C}$ ammette sezioni miste: una persona può ospitare regimi diversi su contesti diversi, e su uno stesso contesto può alternare regimi su sottosequenze distinte. Che $G$ su $\mathcal{I}_{c_{AG}}$ sia stata stabilmente in $\LP$, senza nessuna sottosequenza in $\Lt$, è fatto del caso specifico --- non teorema del formalismo. Altri casi possono avere copertura mista. Il libro non costruisce il framework per spiegare questo caso: il caso è istanza di un framework che ammette molti altri pattern.

\paragraph{Disallineamento simmetrico vs.~asimmetria della copertura.} Due piani da non collassare. Il disallineamento è il teorema di incompatibilità: $\Lt$ e $\LP$ tirano in direzioni opposte su qualunque contesto in cui coabitino. L'asimmetria della copertura è il fatto del caso: dal lato $G$, su questo $c$, non c'è stato $\Lt$ neanche per un momento. Il primo è struttura; il secondo è dato empirico. Servono entrambi per leggere il caso, ma sono cose diverse.

\chapter{Il pianto come errore di tipo}

Lei aveva proposto la sera al bar di Vicenza, dopo gli sms espliciti di $A$. $A$ si alzò e disse la frase semplice: \emph{Greta, non ti chiedo di fidarti di me. Ti dico che vorrei costruire con te quella fiducia.} Lei pianse oltre ogni misura.

\section*{Tipo non formabile in $\GammaG$}

$A$ emette un atto $\alpha$ --- la dichiarazione --- il cui tipo nel sé $\Lt$ è ``richiesta di passaggio di regime su $c$ a $\Lt$'': peso elevato, senza armatura, senza calibrazione strategica. La tradizione chiama questa postura \emph{Patroclo} --- pura apertura.

In $\GammaG$ il tipo non è formabile. $\alpha$ non è processabile come configurazione del pool (troppo carica), né filtrabile come minaccia (non del tipo che $-\lambda \cdot \mathrm{threat}$ scarta). Viene attraversato senza registrazione --- non è esclusione, perché l'esclusione richiederebbe classificazione. È scope error: il tipo non è nel contesto.

La risposta a questa non-tipabilità è somatica, non semantica: pianto. Il pianto è output strutturale di $\LP$ di fronte a un input fuori-grammatica --- non rifiuto, non accoglienza, non strategia. La grammatica non ha l'operazione richiesta; il corpo prende il posto della grammatica che manca.

\section*{Il doppio fraintendimento al livello-grammatica}

Dal lato $\Lt$, $A$ legge il pianto come acknowledgment: ``ha capito la mia richiesta''. Dal lato $\LP$, il pianto non porta significato semantico. I due lati hanno letto lo stesso atto sotto due grammatiche diverse, e nessuno dei due ha mentito.

Fraintendimento di livello-grammatica, anteriore al livello-contenuto. Da scrivere senza colpa, senza accusa, ma senza addolcire: il pianto non è stato sentito come acknowledgment dal lato che lo ha prodotto. Era errore di tipo che si è espresso somaticamente.

\section*{Achille dopo Patroclo}

Quello che resta di $A$ dopo il bar --- l'algidità della sera, l'ariete a Vicenza, la non-reazione al compleanno --- è Achille. Più preciso, più chiuso, più solo. Non per scelta morale: per struttura. La postura di Patroclo --- apertura senza armatura --- era praticabile fino a quando la struttura dell'altra non l'aveva ancora attraversata senza vederla. Dopo, quella postura non è più disponibile.

\section*{Parallelismo delle conservazioni}

Né $A$ né $G$ hanno scelto. $A$ ha emesso $\alpha$ in $\sigma_{\Lt}$ perché è quello che $\Lt$ produce su un contesto con quel peso accumulato. $G$ ha pianto perché è quello che $\LP$ produce di fronte a un input fuori-grammatica. Nessuno dei due ha potuto fare altrimenti.

La sofferenza è stata reale --- intera, fisica, su entrambi i lati. Non imputabile a nessuno come scelta morale. Due strutture conservative si sono incontrate: il loro parallelismo è quello che produce dolore --- non il conflitto. Nella struttura di Noether, nessuna delle due rompe la conservazione dell'altra. Entrambe continuano fedeli alla propria simmetria. La tragedia non è collisione: è parallelismo delle conservazioni.

\chapter{La funzione di questo libro}

Questo libro è prodotto da $\Lt$, vent'anni dopo l'evento. Il regime $\Lt$ dell'autore sul contesto $c_{AG}$ è restato $\Lt$ nel tempo --- perché $\Lt$ accumula, non dimentica --- e ha esercitato a posteriori la capacità di sospensione del default. La sospensione ha rivelato che l'attribuzione di vent'anni prima --- ``$G$ produce sotto $\Lt$'' --- era falso testimone. La verifica algebrica ha mostrato che $G$ produceva sotto $\LP$, e che il falso testimone era strutturalmente inevitabile per chiunque fosse in $\Lt$ in buona fede.

\section*{Ritrattazione formale}

Il libro è ritrattazione formale del falso testimone. Non rivendica, non accusa, non perdona. Corregge un'attribuzione di legge. Atto unilaterale, possibile solo da $\Lt$, e questo va dichiarato apertamente.

Non chiede a $G$ di fare altrettanto. $\LP$ non ha grammatica per farlo --- manca l'aggiunto sinistro che permetterebbe a $G$ di sollevare la propria sequenza nella propria categoria target, riaprirla, classificare un proprio sollevamento precedente come errore di fibra. Né colpa, né dovere mancato. Strutture diverse, capacità diverse.

\section*{Una finestra che non si replica}

L'osservazione formale è stata possibile in una sovrapposizione precisa. Da una parte $A$ era arrivato al punto in cui la struttura era visibile --- non prima, perché il magone alla porta serviva intero; non dopo, perché dopo l'avrebbe saputo nominare ma non l'avrebbe più attraversato. Dall'altra parte $G$ aveva chiuso il jazz bar come l'aveva chiuso --- non prima, perché non aveva ancora costruito la sera; non dopo, perché dopo avrebbe avuto altri argomenti su cui esercitare il copione, e $A$ sarebbe stato un dato vecchio della mappa.

Le due finestre si sono sovrapposte per qualche ora. In quelle ore l'osservazione era possibile. Fuori da quelle ore non lo sarebbe stata. Proprietà della relazione tra le due geometrie, non merito o sfortuna di nessuno.

Il libro è quello che resta della finestra. Non c'è altro da raccogliere; non c'è altro che si possa fare con quei dati che non sia il fissaggio.

\section*{Il bordo del logos}

Il libro ha percorso fino al bordo. Sulle leggi, sulle coordinate, sulle dinamiche, sui tempi, sull'errore di fibra, sulla posizionalità --- su tutto questo, il logos parla.

Oltre il bordo ci sono le domande che il logos non può portare a parola: il valore comparativo delle leggi, la normatività degli incontri tra esse, il significato della sofferenza che producono, il rapporto metafisico tra le due forme di amore, il contenuto del sogno della persona sotto $\LP$. Esempi. Altre domande analoghe --- sul tempo, sull'identità, sul ruolo dell'osservatore --- restano oltre allo stesso modo. Tutte appartengono a un ordine diverso da quello strutturale.

\paragraph{Tre rifiuti, una postura.} La massima distanza spettrale tra $\Lt$ e $\LP$ (cap.~\emph{Lo spettro}) esclude una famiglia continua $L_\alpha$ che le connetta: ogni riduzione di simmetria abbassa il rango o aumenta la torsione, e i due estremi non hanno punto medio raggiungibile da entrambi. Il sogno --- l'orientamento pre-riflessivo che alimenta $\LP$ --- è dato esterno, non derivabile dalla legge. Una $\mathcal{L}^{*}$ unificante richiederebbe una posizione meta che la posizionalità linguistica esclude. Tre rifiuti, statuti diversi --- algebrico, fenomenologico, ontologico --- una sola postura: il logos non chiude.

\paragraph{Il principio si auto-applica.} Il libro è scritto in $\Gamma_{\Lt}$. Dichiara il proprio limite senza uscirne. Non meta-posizione --- la meta-posizione non esiste. È registrazione del proprio errore di tipo come parte dell'oggetto descritto.

\bigskip

Non lo so. Il libro non lo sa. Il logos del libro è onesto abbastanza da non rispondere a queste domande --- non perché siano irrilevanti, ma perché appartengono a un ordine che la dimostrazione non raggiunge. Risposte a questo livello sarebbero proiezioni di una delle due leggi sull'altra, o postulazione di un meta-osservatore che non c'è. Sarebbero più dannose del silenzio.

\bigskip

La tragedia greca finiva con il coro che, dopo aver visto la catastrofe, taceva o pronunciava una formula generale --- non una spiegazione. Il bordo del logos di questo libro è quel silenzio. Il libro ha detto tutto quello che può dire. Quello che resta non è per la prosa.

\bigskip\bigskip
\begin{center}
\emph{Fine.}
\end{center}

% ----------------- APPENDICE -----------------
\part{Appendice: la categoria delle leggi candidate}

\noindent\itshape $\Lt$ e $\LP$ sono due punti su un asse. Il cap.~\emph{Lo spettro} ha nominato gli altri abitanti nella tabella delle coordinate; questa appendice li articola: schema $(S, \Gamma, \rho)$, esempio fenomenologico tratto dalla vita ordinaria, struttura della famiglia parametrica $\Lr^{(n)}$ che interpola tra $\Lt$ e $\LP$. La forma resta sketchata --- ciascuna legge meriterebbe il proprio libro --- ma quanto basta a mostrare che $\Lt$ e $\LP$ non sono speciali, e che la coppia (legge, contesto) è il livello dove sano e patologico si decidono.\upshape

\bigskip

\section*{$\Lo$ --- compulsioni}

$S_{\Lo} = \mathbb{Z}^{|\Sigma|-1} \oplus \mathbb{Z}/n\mathbb{Z}$: monoide libero su atti non rituali $\times$ contatore ciclico sul rituale. $\Gamma_{\Lo} = \mathrm{Aut}(\Sigma) \times \mathbb{Z}/n\mathbb{Z}$. La rottura realizza l'invarianza del deposito sugli atti non rituali e la ciclicità del rituale modulo $n$. Carica: $\Phii$ sui non rituali e posizione ciclica sul rituale. Coordinata: $(|\Sigma|-1,\,\mathbb{Z}/n)$.

Il rituale chiude su sé stesso dopo $n$ esecuzioni e ricomincia. La memoria è cumulativa sui non rituali, periodica sul rituale. La metacognizione è limitata: l'aggiunto sinistro a $U_{\Lo}$ esiste sulla parte libera, manca sulla parte ciclica --- $\Lo$ può sospendere atti, non può sospendere il rituale finché il rituale è dentro il proprio ciclo.

\paragraph{Esempio.} Lo studente che apre il libro, sottolinea il titolo del capitolo, ricopia l'indice in margine, sottolinea di nuovo, ricopia di nuovo --- $n$ volte --- prima di poter leggere la prima riga. Il rituale è $n$-ciclico; gli atti del leggere e capire stanno sulla parte libera. La parte ciclica non ammette sospensione finché non chiude il ciclo; la parte libera sì.

\section*{$\Lr$ --- routine}

$S_{\Lr} = \mathbb{Z}/n\mathbb{Z}$: solo posizione nel ciclo. $\Gamma_{\Lr} = \mathbb{Z}$, traslazione. La rottura realizza la ciclicità modulo $n$. Carica: posizione nel ciclo. Coordinata: $(0,\,\mathbb{Z}/n)$.

\paragraph{Esempio.} Il sonno e la veglia, le tre colazioni della settimana, il ciclo delle stagioni nel calendario domestico. Non c'è accumulo lungo la traiettoria --- ogni mattina è la stessa mattina, indicizzata dalla sua posizione nel ciclo settimanale. Non c'è scelta da fare passo per passo --- la posizione $k+1$ è data dalla posizione $k$. È legge senza deliberazione e senza accumulo.

\paragraph{La famiglia parametrica $\Lr^{(n)}$.} La routine è famiglia indicizzata da $n \in \mathbb{N}$, con due limiti notevoli:

\emph{$n \to \infty$.} Quando il ciclo non si chiude mai, $\mathbb{Z}/n$ tende strutturalmente a $\mathbb{Z}$. La routine smette di essere routine: ogni posizione è distinta da tutte le precedenti, l'asse diventa libero, si recupera la coordinata $(\infty, 0)$ di $\Lt$. Fenomenologicamente: una vita che non torna su sé stessa --- ogni giorno un atto nuovo che si deposita.

\emph{$n \to 1$.} Quando il ciclo collassa a un solo stato, $\mathbb{Z}/n = \{*\}$ e la posizione nel ciclo è banale. Resta solo la configurazione corrente, senza nemmeno la memoria periodica. Si recupera la coordinata $(0, 0)$ di $\LP$. Fenomenologicamente: un presente che non distingue il momento del ciclo, perché il ciclo è uno solo.

$\Lr$ è quindi interpolazione esplicita tra $\Lt$ e $\LP$ sull'asse del rango/torsione. Le due leggi del libro sono i due limiti di una stessa famiglia.

\section*{$\Lf$ --- fobie}

$S_{\Lf} = P$, pool con punto fisso $x \in P$. $\Gamma_{\Lf} = \Sym(P)$. La rottura realizza lo stabilizzatore $\mathrm{Stab}(x) \subset \Sym(P)$. Carica: distanza da $x$. Coordinata: $(0,\,0)$ con vincolo di punto fisso.

Differisce da $\LP$ per la presenza di $x$: una configurazione che il sistema non può attraversare. La dinamica è scelta minima con vincolo. La fobia è il pattern che il vincolo produce sul cammino.

\paragraph{Esempio.} Il claustrofobico che organizza ogni spostamento per non passare in ascensore: $x$ = ``trovarsi in spazio chiuso senza uscita visibile''. La traiettoria del cammino quotidiano evita $x$ per costruzione, e tutto il resto del comportamento --- scelta di scale, di edifici, di percorsi --- è $\mathrm{Stab}(x)$-invariante: permuta tutto quello che non è $x$. La fobia non è il timore dello spazio chiuso; è la forma che il cammino prende sotto il vincolo che $x$ è inaccessibile.

\section*{Sano e patologico come proprietà della coppia}

Non è la legge a essere patologica. È la coppia $(L, c)$. Una legge $L$ governa un contesto $c$ sanamente se $L(c)$ è compatibile senza forzature --- il contesto sopporta naturalmente la struttura. Patologia è compatibilità a forza: la carica si conserva ma il contesto è deformato per sopportarla.

Tre coppie, tre verifiche.

\emph{$\Lo$ sul contesto del lavoro pianistico: sano.} Il contesto della scala ammette rituale --- il rituale \emph{è} la struttura del lavoro. La parte ciclica di $\Lo$ trova nel contesto la propria forma naturale, niente deformazione.

\emph{$\Lo$ sul contesto dell'uscita di casa: patologico.} Il contesto del controllo del gas non ammette rituale --- chiudere la chiave una volta basta, il fatto è acquisito. La parte ciclica di $\Lo$ chiede $n$ verifiche dove il contesto ne ammette una. La coppia regge per forza; il prezzo è il tempo perso e l'ansia residua del ciclo che non chiude.

\emph{$\Lr$ sul contesto del lutto: patologico.} Una routine ben formata serve sul contesto delle stagioni domestiche; sul contesto del lutto chiede al soggetto di trattare il giorno della scomparsa come una posizione nel ciclo settimanale, indistinguibile dalle altre, ricorrente. Il lutto è invece evento unico, non posizione ciclica: il contesto richiede $\Lt$, la sua coordinata libera, la sua dinamica variazionale che pesa la perdita come deposito non saldato. Forzare $\Lr$ sul lutto è patologia: la carica si conserva (ogni giorno è uguale all'altro), ma il dolore non trova la struttura che gli darebbe forma. Diventa rumore di fondo che non si scarica.

Stessa legge, contesti diversi, statuti diversi. La patologia non è proprietà dell'individuo, e nemmeno della legge: è proprietà della coppia. Relazionale, non assoluta.

\end{document}
